\paragraph{Evaluation population.}
The deployments are GPT-5.6 Sol, Claude Opus 5, GPT-5.4, Grok 4.3, Kimi K3, Claude Opus 4.8, DeepSeek V4 Pro. The original-instruction cohort contains
107,520 responses: 15,360 per deployment, with the same 128 held-out
prompts in each of five domains and three levels, and eight stateless responses
per prompt. Requests use no model-side tools or conversation history, and
the evaluation includes no task training or replay intervention.

\paragraph{Provider configurations.}
All deployments request medium reasoning effort and an 8,192-token native
output limit. Claude Opus 5 and Claude Opus 4.8 use Anthropic Messages with
adaptive thinking; the other deployments use hosted API routes.

The requested reasoning setting and output limit do not match compute across
providers: effort labels and token accounting have provider-specific meanings.
Temperature and nucleus sampling use provider defaults, whose exact values
are unknown when the provider does not return them.

Strict grading checks the original executable answer; formatting-normalized
grading applies the same typography corrections across deployments before
the executable checks. The accuracy row of the main display uses normalized
grades and the revised Opus 5 \texttt{Python} instruction at all three levels;
its other accuracy cells use the original instructions. The diversity row
uses strict grades and original instructions for every deployment.
The tables below report original-instruction results under both grading rules.
Prompt preferences and the Level-1 \texttt{PantryPlan} interface difference limit
comparisons across levels. \texttt{Countdown} has no certified finite-total-support
uniform reference. This inference-only comparison does not identify a training
cause, a model-size effect, or a difficulty effect. Collision intervals are
pointwise estimates for individual cells, not intervals for paired deployment differences.

The 15,360 scored DeepSeek V4 Pro responses include verification failures
and truncations. Usage totals omit requests without reported usage and
therefore give a lower bound on the cost of all attempted requests.
